Economic communities of interest---workers in the same industry cluster,
residents sharing a commercial corridor, communities organized around a
manufacturing employer---are among the oldest recognized forms of community
of interest in redistricting law.
Yet most algorithmic redistricting systems have no mechanism to preserve them.
This paper presents M.1, the first paper in the M-track community character
series, which operationalizes economic community membership through the
Census Bureau's LODES Workplace Area Characteristics (WAC) data.

For each census tract we compute a three-dimensional economic character vector:
commercial intensity (fraction of jobs in retail, finance, information, and
real estate), industrial fraction (fraction of jobs in agriculture, mining,
manufacturing, and transportation), and jobs per resident (total jobs divided
as an employment intensity proxy capped at 1.0, keeping all components on $[0,1]$).
Edge weights between adjacent tracts are modulated by the cosine similarity
between their character vectors: similar tracts (both residential, both
commercial, both industrial) receive heavier edges that METIS avoids cutting;
dissimilar tracts receive lighter edges that invite cuts at economic zone
transitions.

Evaluated on North Carolina ($k=14$), Wisconsin ($k=8$), and Texas ($k=38$),
the method produces a 14.3\,pp proportionality gap improvement in North
Carolina (from $-$20.7\,pp to $-$6.4\,pp) as the Research Triangle Park
employment concentration and the Charlotte commercial corridor emerge as
natural cut seams.
Wisconsin shows no effect, as expected when economic geography is diffuse.
Texas shows a 2.6\,pp worsening, consistent with industrial clusters being
embedded in Republican-leaning suburban areas.

The M-track primary metric---mean within-district standard deviation of
\texttt{jobs\_per\_resident}---captures economic community coherence directly.
NC results are confirmed from an actual pipeline run; full 50-state
within-district variance analysis is Phase~2.
The LEHD LODES data is annual, free, publicly available, and contains no
racial or partisan content, making this signal among the most legally
defensible community-of-interest operationalizations available to
redistricting practitioners.
